% ============================================================
% TUTORIAL: Fase 4 — File uploads con S3 presigned URLs
% Proyecto: Sentinel AcademIA
% Compilar: tectonic tutorial-fase4-files.tex
% ============================================================

\documentclass[11pt,a4paper,oneside]{article}

\usepackage[utf8]{inputenc}
\usepackage[T1]{fontenc}
\usepackage[spanish,es-nodecimaldot]{babel}
\usepackage{lmodern}
\usepackage[margin=2.5cm]{geometry}
\usepackage{parskip}
\usepackage{microtype}

\usepackage{xcolor}
\usepackage{colortbl}
\usepackage{array}
\usepackage{booktabs}
\usepackage{amssymb}
\usepackage{pifont}
\usepackage{tcolorbox}
\tcbuselibrary{breakable,skins}

\usepackage{listings}

\lstdefinelanguage{yaml}{
  basicstyle=\ttfamily\small,
  keywordstyle=\color{blue!70!black}\bfseries,
  stringstyle=\color{red!70!black},
  commentstyle=\color{green!50!black}\itshape,
  morekeywords={Type,Properties,Version,Statement,Effect,Principal,Action,Resource},
}

\lstdefinelanguage{bash}{
  basicstyle=\ttfamily\small,
  keywordstyle=\color{blue!70!black}\bfseries,
  stringstyle=\color{red!70!black},
  commentstyle=\color{green!50!black}\itshape,
  morekeywords={aws,npm,cd,ls,mkdir,rm,cp,mv,export,unset,echo,cat,curl,sleep,npx,git},
}

\lstdefinelanguage{json}{
  basicstyle=\ttfamily\small,
  showstringspaces=false,
  breaklines=true,
  literate=
    *{\{}{{{\{}}}{1}
     {\}}{{{\}}}}{1}
     {[}{{{[}}}{1}
     {]}{{{]}}}{1}
     {:}{{{:}}}{1}
     {,}{{{,}}}{1},
  string=[s]{"}{"},
  stringstyle=\color{red!70!black},
  morestring=[b]",
}

\lstdefinelanguage{typescript}{
  basicstyle=\ttfamily\small,
  keywordstyle=\color{blue!70!black}\bfseries,
  stringstyle=\color{red!70!black},
  commentstyle=\color{green!50!black}\itshape,
  morekeywords={import, export, from, const, let, var, function, async, await,
                return, if, else, for, while, true, false, null, undefined,
                new, this, class, extends, interface, type, public, private,
                protected, static, readonly, void, never, any, unknown, string,
                number, boolean, object},
}

\lstset{
  basicstyle=\ttfamily\small,
  numberstyle=\tiny\color{gray},
  numbers=left,
  numbersep=8pt,
  frame=single,
  framerule=0.4pt,
  rulecolor=\color{gray!50},
  backgroundcolor=\color{gray!5},
  breaklines=true,
  breakatwhitespace=true,
  showstringspaces=false,
  captionpos=b,
  tabsize=2,
  literate={á}{{\'a}}1 {é}{{\'e}}1 {í}{{\'i}}1 {ó}{{\'o}}1 {ú}{{\'u}}1
           {ñ}{{\~n}}1 {Á}{{\'A}}1 {É}{{\'E}}1 {Í}{{\'I}}1 {Ó}{{\'O}}1 {Ú}{{\'U}}1
           {¿}{{?`}}1 {¡}{{!`}}1
}

\usepackage[colorlinks=true,linkcolor=blue!60!black,urlcolor=blue!60!black,citecolor=blue!60!black]{hyperref}

\usepackage{fancyhdr}
\pagestyle{fancy}
\fancyhf{}
\fancyhead[L]{\small\textsc{Sentinel AcademIA}}
\fancyhead[R]{\small Tutorial Fase 4}
\fancyfoot[C]{\thepage}
\renewcommand{\headrulewidth}{0.4pt}

\definecolor{okgreen}{RGB}{40,140,60}
\definecolor{errred}{RGB}{200,50,50}
\definecolor{warnorange}{RGB}{220,130,30}
\definecolor{infoblue}{RGB}{30,90,180}

\newcommand{\ok}{\textcolor{okgreen}{\textbf{\checkmark}}}
\newcommand{\err}{\textcolor{errred}{\textbf{\ding{55}}}}
\newcommand{\warn}{\textcolor{warnorange}{\textbf{\textasciitilde}}}

\title{\Huge\textbf{Fase 4: File Uploads} \\[0.3em]
       \Large S3 Presigned URLs + Async Processing \\[0.5em]
       \normalsize Tutorial paso a paso}
\author{Proyecto Sentinel AcademIA \\ \small lFunknown}
\date{\today}

\begin{document}

\maketitle
\thispagestyle{empty}

\begin{tcolorbox}[colback=blue!5!white,colframe=infoblue,title={\textbf{¿Qué construimos en este tutorial?}}]
En la \textbf{Fase 4} permitimos que los usuarios suban \textbf{archivos de evidencia} (fotos, PDFs) a sus quejas. Esto se hace con el patr\'on de \textbf{presigned URLs}:

\begin{enumerate}
  \item Cliente crea queja (POST /api/quejas) -- ya exist\'ia (Fase 1)
  \item Cliente pide URL de upload (POST /api/quejas/\{id\}/upload-url)
  \item API devuelve URL temporal (v\'alida 15 min)
  \item Cliente sube el archivo \textbf{directamente a S3} (sin pasar por la API)
  \item S3 dispara un evento que llama a \texttt{file\_processor} Lambda
  \item Lambda actualiza la queja en DynamoDB con la URL del archivo
\end{enumerate}

\textbf{Decisiones clave}:
\begin{itemize}
  \item \textbf{Por qu\'e presigned URLs}: el cliente sube directo a S3, evitando saturar la API
  \item \textbf{Por qu\'e S3 event}: desacopla el upload del procesamiento
  \item \textbf{Qu\'e se logr\'o}: end-to-end upload funcionando (POST $\rightarrow$ URL $\rightarrow$ S3 PUT $\rightarrow$ Lambda $\rightarrow$ DynamoDB $\rightarrow$ GET con adjuntos)
\end{itemize}
\end{tcolorbox}

\begin{tcolorbox}[colback=warnorange!5!white,colframe=warnorange,title={\textbf{Email (SES) -- no implementado en esta versi\'on}}]
La parte de notificar por email (SES) \textbf{no se implement\'o} en esta versi\'on. El \texttt{analyze\_queja} actualiza el status a \texttt{ANALIZADA} pero no env\'ia email. Para prod: implementar \texttt{notifyByEmail} con AWS SES.
\end{tcolorbox}

\tableofcontents
\newpage

% ============================================================
\section{El patr\'on de presigned URLs}
% ============================================================

\subsection{¿Qu\'e es una presigned URL?}

Una URL temporal que \textbf{S3 genera} y que permite a un cliente \textbf{subir archivos directamente a S3} sin tener credenciales de AWS.

\textbf{Caracter\'isticas}:
\begin{itemize}
  \item \ok \textbf{Temporal}: expira (configurable, default 15 min)
  \item \ok \textbf{Espec\'ifica}: solo PUT (o GET) sobre UN key espec\'ifico
  \item \ok \textbf{Scoped}: puede limitar ContentType, ContentLength
  \item \ok \textbf{No requiere credenciales} del cliente
\end{itemize}

\subsection{¿Por qu\'e este patr\'on?}

\textbf{Alternativa naive}: el cliente env\'ia el archivo a la API, la API lo sube a S3.

\begin{itemize}
  \item \err \textbf{Problema 1}: la API recibe archivos grandes (10MB+), la Lambda tiene l\'imite de 6MB de payload (sincr\'ono) o 250MB (async)
  \item \err \textbf{Problema 2}: desperdicia bandwidth y CPU de la Lambda
  \item \err \textbf{Problema 3}: timeout si el upload es lento
  \item \ok \textbf{Soluci\'on con presigned URL}: la API solo genera la URL, el cliente sube directo
\end{itemize}

\subsection{Diagrama del flujo}

\begin{center}
\begin{tabular}{ccccccc}
\textbf{Cliente} & $\rightarrow$ & \textbf{API} & $\rightarrow$ & \textbf{S3 presigned} \\
& & genera URL & & PUT directo \\
\end{tabular}
\end{center}

\begin{center}
\begin{tabular}{ccc}
\textbf{S3} & $\rightarrow$ \textbf{file\_processor} & $\rightarrow$ \textbf{DynamoDB} \\
evento s3:ObjectCreated & actualiza queja & guarda URL \\
\end{tabular}
\end{center}

% ============================================================
\section{Componente 1: S3 Bucket}
% ============================================================

\subsection{Crear el bucket en SAM}

\begin{lstlisting}[language=yaml,caption={infra/template.yaml — S3 bucket}]
FilesBucket:
  Type: AWS::S3::Bucket
  Properties:
    BucketName: !Sub "sentinel-files-${Stage}-${AWS::AccountId}"
    PublicAccessBlockConfiguration:
      BlockPublicAcls: true
      BlockPublicPolicy: true
      IgnorePublicAcls: true
      RestrictPublicBuckets: true
    CorsConfiguration:
      CorsRules:
        - AllowedOrigins: ["*"]
          AllowedMethods: [GET, PUT, POST]
          AllowedHeaders: ["*"]
          MaxAge: 3000
\end{lstlisting}

\textbf{Decisiones}:
\begin{itemize}
  \item \ok \textbf{Nombre \'unico}: incluye account ID para evitar colisiones
  \item \ok \textbf{PublicAccessBlock: todo ON}: el bucket NO es p\'ublico
  \item \ok \textbf{Acceso v\'ia presigned URL}: solo usuarios autorizados
  \item \ok \textbf{CORS}: permite PUT desde el frontend
\end{itemize}

% ============================================================
\section{Componente 2: file\_service.py (presigned URLs)}
% ============================================================

\subsection{Funci\'on principal: \texttt{generate\_presigned\_upload\_url}}

\begin{lstlisting}[language=Python,caption={src/services/file\_service.py}]
import boto3
import re
import uuid

ALLOWED_CONTENT_TYPES = [
    "image/jpeg", "image/png", "image/gif", "image/webp",
    "application/pdf",
    "application/msword",
    "application/vnd.openxmlformats-officedocument.wordprocessingml.document",
    "text/plain",
]
MAX_FILE_SIZE_BYTES = 10 * 1024 * 1024  # 10 MB

def build_s3_key(tenant_id, queja_id, filename):
    """Pattern: tenants/{tenantId}/quejas/{quejaId}/{uuid}-{filename}"""
    safe_filename = re.sub(r"[^a-zA-Z0-9._-]", "_", filename)
    file_id = uuid.uuid4().hex[:12]
    return f"tenants/{tenant_id}/quejas/{queja_id}/{file_id}-{safe_filename}"


def generate_presigned_upload_url(
    queja_id, tenant_id, filename, content_type, file_size=None
):
    settings = get_settings()
    bucket = settings.files_bucket_resolved

    # Validar content type
    if content_type not in ALLOWED_CONTENT_TYPES:
        raise FileError(f"Content type {content_type} not allowed")

    # Validar tama\~no (si se proporcion\'o)
    if file_size is not None and file_size > MAX_FILE_SIZE_BYTES:
        raise FileError(f"File too large: {file_size} bytes (max {MAX_FILE_SIZE_BYTES})")

    s3_key = build_s3_key(tenant_id, queja_id, filename)
    s3 = boto3.client("s3")
    expires_in = 900  # 15 minutos

    upload_url = s3.generate_presigned_url(
        "put_object",
        Params={
            "Bucket": bucket,
            "Key": s3_key,
            "ContentType": content_type,
        },
        ExpiresIn=expires_in,
    )

    return {
        "uploadUrl": upload_url,
        "fileKey": s3_key,
        "fileUrl": build_s3_url(bucket, s3_key),
        "expiresIn": expires_in,
        "maxSize": MAX_FILE_SIZE_BYTES,
    }
\end{lstlisting}

\textbf{Decisiones}:
\begin{itemize}
  \item \ok \textbf{Key pattern}: \texttt{tenants/\{tenantId\}/quejas/\{quejaId\}/...} (multi-tenant)
  \item \ok \textbf{Sanitize filename}: solo alfanum\'erico, puntos, gui\'on
  \item \ok \textbf{Validar ContentType}: solo PDFs, im\'agenes, docs
  \item \ok \textbf{Validar tama\~no}: 10MB max
  \item \ok \textbf{Expiraci\'on}: 15 min (suficiente para upload)
\end{itemize}

% ============================================================
\section{Componente 3: upload\_url Lambda}
% ============================================================

\subsection{El handler}

\begin{lstlisting}[language=Python,caption={src/handlers/upload\_url.py}]
from pydantic import BaseModel, Field

class UploadURLRequest(BaseModel):
    filename: str = Field(min_length=1, max_length=255)
    contentType: str = Field(min_length=1, max_length=100)
    fileSize: int | None = Field(default=None, ge=0, le=10_485_760)

@error_handler
def lambda_handler(event, context):
    with_correlation_id(event)
    tenant_id = extract_tenant_id(event)
    log.append_keys(tenant_id=tenant_id)

    # Extraer quejaId del path
    path_params = event.get("pathParameters") or {}
    queja_id = path_params.get("quejaId") or path_params.get("id")
    if not queja_id:
        raise ValidationError(message="Missing quejaId in path")

    # Verificar que la queja existe (multi-tenant safety)
    dynamo = get_dynamo_client()
    dynamo.get_queja(tenant_id=tenant_id, queja_id=queja_id)  # raises 404

    # Parse + validate body
    body = _parse_body(event)
    req = UploadURLRequest.model_validate(body)

    # Generar presigned URL
    result = generate_presigned_upload_url(
        queja_id=queja_id,
        tenant_id=tenant_id,
        filename=req.filename,
        content_type=req.contentType,
        file_size=req.fileSize,
    )

    return HttpResponse.ok(result)
\end{lstlisting}

\textbf{Patrones importantes}:
\begin{itemize}
  \item \ok \textbf{Verifica que la queja existe} antes de generar URL (404 si no)
  \item \ok \textbf{Valida con Pydantic} (ContentType, FileSize, etc.)
  \item \ok \textbf{tenantId del header, no del body} (seguridad)
\end{itemize}

% ============================================================
\section{Componente 4: file\_processor Lambda (S3 event consumer)}
% ============================================================

\subsection{El handler}

\begin{lstlisting}[language=Python,caption={src/handlers/file\_processor.py}]
import re
from typing import Any

def _parse_key(s3_key):
    """Extrae tenant_id y queja_id del S3 key."""
    match = re.match(r"^tenants/([^/]+)/quejas/([^/]+)/", s3_key)
    if not match:
        return None
    return match.group(1), match.group(2)


def _process_record(record):
    """Procesa un solo record de S3."""
    bucket = record["s3"]["bucket"]["name"]
    s3_key = record["s3"]["object"]["key"]

    parsed = _parse_key(s3_key)
    if not parsed:
        log.warning("Could not parse S3 key")
        return
    tenant_id, queja_id = parsed

    # Construir URL p\'ublica del archivo
    file_url = build_s3_url(bucket, s3_key)

    # Actualizar la queja con esta URL (append a adjuntos)
    dynamo = get_dynamo_client()
    item = dynamo.get_queja(tenant_id=tenant_id, queja_id=queja_id)
    adjuntos = item.get("adjuntos", []) or []
    if file_url not in adjuntos:
        adjuntos.append(file_url)
        dynamo.update_item(
            pk=queja_pk(tenant_id, queja_id),
            sk=queja_sk(queja_id),
            updates={"adjuntos": adjuntos, "updatedAt": now_iso()},
        )


@error_handler
def lambda_handler(event, context):
    for record in event.get("Records", []):
        try:
            _process_record(record)
        except Exception:
            log.exception("Failed to process S3 record")
    return {"statusCode": 200, "body": "OK"}
\end{lstlisting}

\textbf{Patrones}:
\begin{itemize}
  \item \ok \textbf{Parse S3 key}: \texttt{tenants/\{tenantId\}/quejas/\{quejaId\}/...} $\rightarrow$ extrae ambos IDs
  \item \ok \textbf{Append, not overwrite}: si el archivo ya est\'a, no duplica
  \item \ok \textbf{Continue on error}: si un record falla, sigue con los dem\'as
\end{itemize}

% ============================================================
\section{Componente 5: SAM template (los Lambdas nuevos)}
% ============================================================

\subsection{Upload URL Lambda + permisos}

\begin{lstlisting}[language=yaml,caption={infra/template.yaml — UploadURLFunction}]
UploadURLFunction:
  Type: AWS::Serverless::Function
  Properties:
    Role: !Ref LambdaRoleArn
    FunctionName: !Sub "sentinel-upload-url-${Stage}"
    CodeUri: ../src
    Handler: handlers.upload_url.lambda_handler
    Environment:
      Variables:
        DYNAMODB_TABLE: !Sub "sentinel-quejas-${Stage}"
        FILES_BUCKET: !Ref FilesBucket
    Events:
      UploadApi:
        Type: Api
        Properties:
          Path: /api/quejas/{quejaId}/upload-url
          Method: POST
          RestApiId: !Ref SentinelApi
    Policies:
      - S3ReadPolicy:
          BucketName: !Ref FilesBucket
      - Statement:
          - Effect: Allow
            Action: s3:PutObject
            Resource: !Sub "arn:aws:s3:::${FilesBucket}/*"
\end{lstlisting}

\textbf{Permisos necesarios}:
\begin{itemize}
  \item \ok \texttt{S3ReadPolicy}: para verificar si el objeto existe (opcional)
  \item \ok \texttt{s3:PutObject}: para generar presigned URL (necesario)
  \item \err \textbf{No \texttt{s3:*}}: principio de menor privilegio
\end{itemize}

% ============================================================
\section{El problema: circular dependency}
% =========================================================%

\subsection{El bug que enfrentamos}

Al intentar agregar el S3 event al Lambda, obtuvimos:

\begin{lstlisting}
Error: Circular dependency between resources:
  [SentinelApi, FileProcessorFunction, CreateQuejaFunctionCreateApiPermissionStage,
   FilesBucket, ListQuejasFunctionListApiPermissionStage, UploadURLFunction,
   FileProcessorFunctionS3EventPermission, ...]
\end{lstlisting}

\subsection{¿Por qu\'e?}

El problema:
\begin{itemize}
  \item \textbf{FilesBucket} (con NotificationConfiguration) \textbf{necesita} el ARN de FileProcessorFunction
  \item \textbf{FileProcessorFunction} \textbf{depende} de FilesBucket (para permisos)
  \item $\rightarrow$ \textbf{dependencia circular}
\end{itemize}

\subsection{La soluci\'on: configurar S3 notification manual}

En vez de usar SAM event source (que crea el notification autom\'aticamente), definimos el Lambda Permission por separado y configuramos la notificaci\'on con AWS CLI post-deploy.

\begin{lstlisting}[language=yaml,caption={infra/template.yaml — permission separada}]
FileProcessorFunction:
  Type: AWS::Serverless::Function
  Properties:
    Role: !Ref LambdaRoleArn
    FunctionName: !Sub "sentinel-file-processor-${Stage}"
    CodeUri: ../src
    Handler: handlers.file_processor.lambda_handler
    Environment:
      Variables:
        DYNAMODB_TABLE: !Sub "sentinel-quejas-${Stage}"
        FILES_BUCKET: !Ref FilesBucket
    # SIN S3 event aqui (causa circular dep)
    Policies:
      - S3ReadPolicy:
          BucketName: !Ref FilesBucket
      - DynamoDBCrudPolicy:
          TableName: !Ref QuejasTable

FileProcessorS3Permission:
  Type: AWS::Lambda::Permission
  Properties:
    FunctionName: !Ref FileProcessorFunction.Arn
    Action: lambda:InvokeFunction
    Principal: s3.amazonaws.com
    SourceAccount: !Ref AWS::AccountId
    SourceArn: !Sub "arn:aws:s3:::${FilesBucket}"
\end{lstlisting}

\begin{lstlisting}[language=bash,caption={Post-deploy: configurar S3 notification}]
BUCKET=$(aws cloudformation describe-stacks \
  --stack-name sentinel-academia-dev \
  --query 'Stacks[0].Outputs[?OutputKey==`FilesBucketName`].OutputValue' \
  --output text)

LAMBDA_ARN=$(aws lambda get-function \
  --function-name sentinel-file-processor-dev \
  --query 'Configuration.FunctionArn' --output text)

aws s3api put-bucket-notification-configuration \
  --bucket "$BUCKET" \
  --notification-configuration "{
    \"LambdaFunctionConfigurations\": [{
      \"Id\": \"FileProcessorTrigger\",
      \"LambdaFunctionArn\": \"$LAMBDA_ARN\",
      \"Events\": [\"s3:ObjectCreated:*\"],
      \"Filter\": {
        \"Key\": {
          \"FilterRules\": [{\"Name\": \"prefix\", \"Value\": \"tenants/\"}]
        }
      }
    }]
  }"
\end{lstlisting}

% ============================================================
\section{Los 2 bugs que arreglamos}
% ============================================================

\begin{tcolorbox}[colback=errred!5!white,colframe=errred,breakable,title={Bug 1: \texttt{from src.handlers.\_shared} en file\_service.py}]
\textbf{Cuando}: upload\_url Lambda retorna 500. \\
\textbf{Causa}: \texttt{services/file\_service.py} tiene imports con \texttt{src.} prefix que ya removimos en Fase 1. \\
\textbf{Soluci\'on}: \texttt{sed -i 's/from src\.handlers/from handlers/g'}.

\begin{lstlisting}
# MAL
from src.handlers._shared.config import get_settings

# BIEN
from handlers._shared.config import get_settings
\end{lstlisting}
\end{tcolorbox}

\begin{tcolorbox}[colback=errred!5!white,colframe=errred,breakable,title={Bug 2: KeyError 'filename' en logger}]
\textbf{Cuando}: log.info() falla con \texttt{KeyError: "Attempt to overwrite 'filename' in LogRecord"}. \\
\textbf{Causa}: \texttt{filename} es palabra reservada del stdlib logging. \\
\textbf{Soluci\'on}: usar \texttt{file\_name} en \texttt{extra=\{\}}.

\begin{lstlisting}
# MAL
log.info("Generated", extra={"filename": req.filename})

# BIEN
log.info("Generated", extra={"file_name": req.filename})
\end{lstlisting}
\end{tcolorbox}

% ============================================================
\section{El test E2E (paso a paso con outputs)}
% =========================================================%

\subsection{Paso 1: Crear queja}

\begin{lstlisting}[language=bash]
curl -s -X POST "$API/api/quejas" \
  -H "Content-Type: application/json" \
  -H "X-Tenant-ID: demo-utpl" \
  -d '{"titulo":"Prueba con archivo adjunto",
       "descripcion":"Subimos evidencia via presigned URL",
       "categoriaDeclarada":"INFRAESTRUCTURA",
       "anonima":false}'
\end{lstlisting}

\begin{lstlisting}[language=json,caption={Output}]
{
  "quejaId": "eb00100c-db86-4a04-8aab-67f3ba357d24",
  "status": "PENDIENTE",
  "correlationId": "...",
  "createdAt": "2026-06-20T14:45:15..."
}
\end{lstlisting}

\subsection{Paso 2: Pedir upload URL}

\begin{lstlisting}[language=bash]
curl -s -X POST "$API/api/quejas/eb00100c-.../upload-url" \
  -H "Content-Type: application/json" \
  -H "X-Tenant-ID: demo-utpl" \
  -d '{"filename":"test.pdf","contentType":"application/pdf","fileSize":1024}'
\end{lstlisting}

\begin{lstlisting}[language=json,caption={Output}]
{
  "uploadUrl": "https://sentinel-files-dev-...s3.amazonaws.com/
    tenants/demo-utpl/quejas/eb00100c-.../c69af945dba7-test.pdf
    ?AWSAccessKeyId=...&Signature=...&Expires=1781967798",
  "fileKey": "tenants/demo-utpl/quejas/eb00100c-.../c69af945dba7-test.pdf",
  "fileUrl": "https://sentinel-files-dev-...s3.us-east-1.amazonaws.com/
    tenants/demo-utpl/quejas/eb00100c-.../c69af945dba7-test.pdf",
  "expiresIn": 900,
  "maxSize": 10485760
}
\end{lstlisting}

\subsection{Paso 3: Subir archivo a S3 (directo, sin API)}

\begin{lstlisting}[language=bash]
echo "%PDF-1.4 fake" > /tmp/test.pdf
curl -s -X PUT "$UPLOAD_URL" \
  -H "Content-Type: application/pdf" \
  --data-binary @/tmp/test.pdf \
  -w "HTTP %{http_code}\n" -o /dev/null
# Esperado: HTTP 200
\end{lstlisting}

\subsection{Paso 4: S3 event $\rightarrow$ file\_processor $\rightarrow$ DynamoDB}

\begin{lstlisting}[language=bash]
sleep 5
curl -s "$API/api/quejas/eb00100c-..." -H "X-Tenant-ID: demo-utpl"
\end{lstlisting}

\begin{lstlisting}[language=json,caption={Output esperado}]
{
  "quejaId": "eb00100c-...",
  "status": "ANALIZADA",
  "adjuntos": [
    "https://sentinel-files-dev-...s3.us-east-1.amazonaws.com/
     tenants/demo-utpl/quejas/eb00100c-.../c69af945dba7-test.pdf"
  ]
}
\end{lstlisting}

% ============================================================
\section{Verificaci\'on en AWS Console}
% =========================================================%

\subsection{Recursos desplegados}

\begin{lstlisting}[language=bash]
aws cloudformation describe-stack-resources \
  --stack-name sentinel-academia-dev \
  --query 'StackResources[?contains(ResourceType, `S3`) || contains(ResourceType, `Lambda`)].{Type:ResourceType, Name:LogicalResourceId, Status:ResourceStatus}' \
  --output table
\end{lstlisting}

\subsection{Ver logs del file\_processor}

\begin{lstlisting}[language=bash]
aws logs tail /aws/lambda/sentinel-file-processor-dev --follow
\end{lstlisting}

Output esperado:
\begin{lstlisting}[language=json]
{"level":"INFO","location":"lambda_handler:65",
 "message":"Processing 1 S3 records"}

{"level":"INFO","location":"_process_record:43",
 "message":"File uploaded: tenants/.../c69af945dba7-test.pdf"}

{"level":"INFO","location":"_process_record:55",
 "message":"Queja updated with new file URL"}
\end{lstlisting}

% ============================================================
\section{Cat\'alogo de errores y soluciones}
% =========================================================%

\begin{tcolorbox}[colback=errred!5!white,colframe=errred,breakable,title={Error 1: Circular dependency}]
\textbf{Cuando}: \texttt{sam deploy} con S3 event source. \\
\textbf{Causa}: bucket y lambda se referencian mutuamente. \\
\textbf{Soluci\'on}: configurar notification con AWS CLI post-deploy.
\end{tcolorbox}

\begin{tcolorbox}[colback=errred!5!white,colframe=errred,breakable,title={Error 2: No module named 'src'}]
\textbf{Cuando}: \texttt{from src.handlers.\_shared.X} en servicios. \\
\textbf{Causa}: en Fase 1 quitamos el prefijo \texttt{src.}, pero services no se actualiz\'o. \\
\textbf{Soluci\'on}: sed para todos los archivos en \texttt{src/services/}.
\end{tcolorbox}

\begin{tcolorbox}[colback=errred!5!white,colframe=errred,breakable,title={Error 3: KeyError 'filename' en logger}]
\textbf{Cuando}: \texttt{log.info(extra=\{``filename'': ...\})}. \\
\textbf{Causa}: \texttt{filename} es palabra reservada en stdlib logging. \\
\textbf{Soluci\'on}: usar \texttt{file\_name} o \texttt{error\_message}.
\end{tcolorbox}

\begin{tcolorbox}[colback=errred!5!white,colframe=errred,breakable,title={Error 4: 403 al hacer PUT al presigned URL}]
\textbf{Cuando}: el cliente intenta subir y recibe 403. \\
\textbf{Causa}: la URL expir\'o o el ContentType no coincide. \\
\textbf{Soluci\'on}: pedir nueva URL con mismo ContentType. \\
\textbf{Para prod}: regenerar URL autom\'aticamente en el frontend.
\end{tcolorbox}

\begin{tcolorbox}[colback=errred!5!white,colframe=errred,breakable,title={Error 5: file\_processor no actualiza la queja}]
\textbf{Cuando}: el upload a S3 funciona pero la queja sigue con \texttt{adjuntos: []}. \\
\textbf{Causa}: el Lambda no se est\'a disparando (S3 notification mal configurada). \\
\textbf{Soluci\'on}: verificar con \texttt{aws s3api get-bucket-notification-configuration}.
\end{tcolorbox}

% ============================================================
\section{Lo que aprendimos (resumen ejecutivo)}
% ============================================================

\begin{tcolorbox}[colback=okgreen!5!white,colframe=okgreen,title={\textbf{Checklist Fase 4}}]
\begin{enumerate}
  \item \textbf{Presigned URLs}: cliente sube directo a S3, no a la API
  \item \textbf{Expiraci\'on}: 15 min (suficiente para upload)
  \item \textbf{Key pattern multi-tenant}: \texttt{tenants/\{tid\}/quejas/\{jid\}/...}
  \item \textbf{ContentType whitelisting}: solo PDFs, im\'agenes, docs
  \item \textbf{Max size 10MB}: validado en backend antes de generar URL
  \item \textbf{S3 event $\rightarrow$ file\_processor}: desacopla upload de update
  \item \textbf{Circular dep}: usar permission manual + post-deploy CLI
  \item \textbf{\texttt{file\_name} en logger}: \texttt{filename} es palabra reservada
  \item \textbf{Sanitize filename}: solo alfanum\'erico, \texttt{.}, \texttt{\_}, \texttt{-}
  \item \textbf{Append, not overwrite}: si el archivo ya est\'a, no duplica
\end{enumerate}
\end{tcolorbox}

% ============================================================
\section{Ejercicios para practicar}
% ============================================================

\subsection{Ejercicio 1: GET /api/quejas/\{id\}/files}

Endpoint que lista los archivos adjuntos de una queja. Usa \texttt{list\_queja\_files} que ya existe en file\_service.

\subsection{Ejercicio 2: DELETE /api/quejas/\{id\}/files/\{key\}}

Eliminar un archivo de S3 y de los adjuntos. Requiere verificar permisos.

\subsection{Ejercicio 3: Image thumbnails}

Si contentType es \texttt{image/*}, generar thumbnail y guardarlo en otro key. Usa \texttt{Pillow} (pero ojo con Lambda size limit).

\subsection{Ejercicio 4: Email notification (Fase 4 parte 2)}

Cuando file\_processor termina, enviar email v\'ia SES al admin del tenant. Requiere SES sandbox configurado.

\subsection{Ejercicio 5: Antivirus scan}

Agregar un Lambda que escanea archivos subidos con ClamAV antes de marcarlos como adjuntos. Patr\'on: S3 event $\rightarrow$ scan $\rightarrow$ tag $\rightarrow$ notify.

% ============================================================
\section{Siguiente paso: Fase 5 (RAG)}
% ============================================================

En la \textbf{Fase 5} agregamos RAG (Retrieval Augmented Generation):

\begin{itemize}
  \item Documentaci\'on del reglamento universitario
  \item Embeddings con Cohere \texttt{embed-multilingual-v3.0}
  \item Vector store (in-memory o FAISS)
  \item Lambda \texttt{chat} que toma una pregunta, busca docs relevantes, llama Cohere con contexto
\end{itemize}

\textbf{Tiempo estimado}: 2-3 horas.

\begin{center}
\begin{tabular}{ll}
\toprule
Fase & Estado \\
\midrule
0. Setup + OCI & [OK] \\
1. Backend Core & [OK] \\
2. LLM integration & [OK] \\
3. Frontend deploy & [OK] \\
4. Files (este tutorial) & [OK] \\
5. RAG & pr\'oximo \\
6. Testing & pendiente \\
7. Deploy final + demo & pendiente \\
\bottomrule
\end{tabular}
\end{center}

\vspace{2em}
\begin{center}
\textit{Fin del tutorial. \`A por la Fase 5!}
\end{center}

\end{document}
